% 本文件由 LaTeX 排版 Agent 根据有效 translation.json 自动生成。
% author=John Chrysostom; work_id=1906; title=论圣巴比拉

\chapter{论圣巴比拉}

% structure: p0001 source=h1 target=omitted reason=page-title-already-rendered-by-parent

1. 我今日急于偿还我近日在此所许下的债务。但我该怎么做呢？与此同时，真福巴比拉显现了，他以面容的光辉吸引我们的注意，虽未发声，却召我到他那里。因此，请不要因我的迟延而不悦；无论如何，时间越长，利息越多。因为我们将把这笔钱连本带利存起来。\BibleRef{路 19:23} 既然那托付给我们的主人这样吩咐了。因此，我深信所借出的本金和利息都在等候你们，让我们不要放过今日遇到的机会，而要在真福巴比拉的英勇事迹中欢欣。\par

他如何治理我们中间的教会，在风暴、波浪和巨浪中拯救那神圣的船；他如何以大胆无畏面对皇帝，如何为羊群舍命，经历那蒙福的屠杀——这些以及诸如此类的事，我们留给年长的导师们和我们的共同父亲去讲述。因为更久远的事，年长者可以向你们讲述；但那些近期发生、在我们有生之年的事，即殉道者死后、埋葬之后、他停留在城郊期间所发生的事，我作为年轻人将向你们讲述。我确实知道希腊人会嘲笑我的承诺，如果我承诺讲述一个已埋葬、已化为尘土的人的死后和埋葬后的英勇事迹。我们当然不会因此保持沉默，反而要特别讲述，为的是通过展示这一奇迹的真相，使他们的嘲笑转而落到他们自己头上。因为对于一个普通人，死后没有英勇事迹；但对于一位殉道者，却有许多伟大的事迹，不是为了使他更显赫（因为他不需要群众的荣耀），而是为了让你——不信者——明白，殉道者的死亡不是死亡，而是更美好生命的开始、更属灵交往的前奏、从坏到好的转变。因此，不要只看到殉道者的身体缺乏灵魂的活力，而要看到有更大的力量代替灵魂临到它身边——我指的是圣神的恩宠，它通过所行的奇迹为所有人辩护复活之事。因为如果天主赐予死去并化为尘土的身体比所有活人更大的能力，那么在他赐予冠冕的时刻，他岂不更要赐予他们比先前更美好、更长久的生活吗？那么这位圣人的英勇事迹是什么呢？但不要烦扰，如果我稍微追溯一下我们的讲述。因为那些想出色展示肖像画的人，不会在把观众稍稍从画前移开之前就揭开画布，而是通过距离使视野更清晰。那么，也请你们对我耐心，让我把叙述引向过去。\par

因为当尤利安——在亵渎神明方面超越所有人——登上皇位，掌握专制权杖时，他立刻举起手来攻击那创造他的天主，无视他的恩主，从下界仰望上天，像疯狗一样狂吠，它们既向不喂养它们的人也向喂养它们的人吠叫。但他比它们更疯狂。因为它们确实对朋友和陌生人同样憎恨；但此人却向与他的救恩无关的魔鬼谄媚，用各种敬拜方式崇拜它们。而对于他的恩主、救主、那为了他而不惜舍弃独生子的那一位，他却转身憎恨，并大肆破坏十字架——那曾扶起匍匐的全世界、驱散四围的黑暗、带来比阳光更灿烂的光明的十字架；即便如此，他仍未停止疯狂，反而许诺要将加里肋亚人的民族从世界中铲除（他惯常如此称呼我们）。然而，如果他认为基督徒的名字是可憎的，基督教本身充满极大羞耻，他为何不以此羞辱我们，却用一个奇怪的名字呢？因为他清楚地知道，被称为属于基督的，不仅对人类，而且对天使和天上的势力，都是极大的装饰。为此，他调动一切，要剥夺我们这装饰，并阻止它的宣讲。但这是不可能的，可怜又可悲的人！正如不可能毁灭天空、熄灭太阳、动摇并推翻大地根基一样，基督早已预言，如此说：\BibleQuote{天地要过去，但是我的话决不会过去。}\BibleRef{玛 24:35}\par

好吧，你不服从基督的话语；那就接受这些行为所发出的宣告吧。因为我确实有幸知道天主的宣告是何等有力、何等不可战胜，我深信它比自然秩序和一切事物的经验更可靠。但你仍匍匐在地，被人类理性的探究所困扰，那么请接受行为的见证。我毫不反对，也不争辩。\par

2. 那么，行为说了什么？基督说过，天地被毁，也比祂的话落空更容易。\BibleRef{路 16:17} 皇帝却反对这些话，并威胁要废除祂的法令。那么，那个威胁这些事的皇帝在哪里？他已毁灭、朽坏，如今在阴府，等待那不可避免的惩罚。但颁布这些法令的基督在哪里？在天堂，坐在圣父的右边，占据荣耀的最高宝座；皇帝的亵渎之言和他那不受管束的舌头在哪里？它们已化为灰烬、尘土和虫子的食物。基督的判决又在哪里？它借着行为的真实发出光芒，从事件的结果获得光辉，如同从金柱上获得光亮一般。然而，当皇帝准备向我们发动战争时，他没有留下任何未做的事，而是召集先知，召唤巫士，到处充满了魔鬼和恶神。\par

那么，这种敬拜的回报是什么？城市的倾覆，最严重的饥荒。因为你们无疑知道并记得，市场是多么空无一物，作坊里充满混乱，人人都争抢最先到手的东西并离开。我何必提及饥荒，就连水泉本身也干涸了，那些水泉曾因水流丰沛而使河流黯然失色。既然我提到了水泉，来，让我们立刻上到\Place{达芙涅}，将我们的论述引向殉道者的高贵事迹。虽然你们仍希望我展示希腊人的无耻，我也希望如此，但让我们避免吧；因为凡是有殉道者纪念的地方，也必有希腊人的羞耻。这位皇帝上到\Place{达芙涅}，常常烦扰阿波罗，祈祷、恳求、哀求，好让未来的事向他预告。那么，这位先知，希腊人的大神说了什么？\Quote{死者阻碍我开口，}他说，\Quote{但要打破坟墓，挖出骨头，移走死者。}还有什么比这些命令更不虔敬的呢？盗墓的魔鬼引入奇怪的法律，设计驱逐陌生人的新方法。谁曾听说过死者被赶走？谁曾见过无生命的尸体被命令移动，如同他所吩咐的，颠覆了自然界的共同规律？因为自然的规律对所有人都是共同的：离世者应被埋在地下，交付埋葬，被覆盖在大地——万物之母的怀抱中；这些规律，无论是希腊人、野蛮人、斯基泰人，或者有比他们更野蛮的人，都从未改变，反而人人尊敬它们，遵守它们，因此它们对所有人都是神圣而可敬的。但魔鬼脱下假面具，光着头抵抗自然界的共同规律。因为他说，死者是一种污染。死者并非污染，最邪恶的魔鬼，邪恶的意念才是可憎的。但如果必须说出惊人之处，充满邪恶的活人身体比死人的身体更污秽。因为活人身体听从心灵的指挥，而死者身体则不动。不动且毫无知觉的东西应免于一切控告。我并非说活人的身体本质上是污秽的；而是说，邪恶乖僻的意念在任何地方都招致全体的指控。\par

那么，尸体并非污染，阿波罗啊，但迫害一位渴望贞洁的少女，侮辱童贞女的尊严，因无耻行为的失败而哀叹，这才是该受指控和惩罚的。无论如何，在我们当中曾有众多奇妙伟大的先知，他们也说了许多关于未来的事，他们从不吩咐问他们的人去挖出死者的骨头。是的，\Person{厄则克耳}站在骸骨本身旁边，不仅没有被它们阻碍，反而给它们加上肉、神经和皮肤，使它们复活。\EditorialNote{厄则克耳第三十七章} 但伟大的\Person{梅瑟}并未站在死者骸骨旁边，而是带走了\Person{若瑟}的整个尸体，从而预言了未来的事。\BibleRef{出 13:19} 这很合理，因为他们的言语是圣神的恩宠。但这些人的言语是欺骗和谎言，绝无法隐藏。因为这些事情是借口和托词，而且皇帝害怕真福\Person{巴比拉}，这从他做的事可以看出来。因为他抛下其他所有死者，只移走了那位殉道者。然而，如果他做这些事是出于厌恶他，而不是出于恐惧，那么他就应当命令将棺材打破，扔进海里，运到沙漠，或用其他毁灭方法使之消失；因为这才是厌恶者的做法。当日天主对希伯来人谈论外邦人的可憎之物时，就是这样做的。他吩咐打碎他们的雕像，而不是将他们的可憎之物从郊外带到城里。\par

3. 于是，殉道者被移动了，但邪神即使在那时也没有摆脱恐惧，反而立刻明白：移动殉道者的骨骸是可能的，但逃脱他的手却是不可能的。因为棺椁一被拉进城，一道雷电就从天而降，击打在那偶像的头上，将其全部烧毁。然而，即便从前不会，至少那时不虔敬的皇帝很可能会发怒，并将他的愤怒发泄在殉道者的见证上。但即使在那时，他也不敢，因为极大的恐惧占据了他。虽然他看到那场焚烧难以忍受，并且确切地知道原因，他却保持沉默。这不仅令人惊奇的是他没有毁掉那见证，而且他甚至不敢再次为那神庙盖上屋顶。因为他知道，他知道，那打击是天主所降下的，他害怕若要制定任何进一步计划，就会招致那火落到自己头上。为此，他忍受看到阿波罗的神龛变得如此荒凉；因为没有别的原因使他没有纠正所发生的事，唯有恐惧。因此，他不情愿地保持沉默，并且明知如此，这给邪神留下了同样多的羞辱，给殉道者留下了同样多的尊荣。因为现在墙壁矗立在那里，代替了战利品，发出比号角更清晰的声音。对\Place{达芙妮}的人、对城里的人、对从远方来的人、对我们这里的人、对以后的人，它们用外观宣布一切：殉道者的角力、斗争、胜利。因为很可能，住在远郊的人，当他看到圣人的小堂失去了神龛，而\Person{阿波罗}的神庙失去了屋顶，就会问及这些事的缘由；然后在了解了全部历史后，就离开那里。这就是殉道者死后所做的崇高事迹，因此，我认为你们的城市是有福的，因为你们对这位圣者表现了极大的热忱。因为那时，当他从\Place{达芙妮}回来时，我们全城的人都涌到路上，市场空无人烟，房屋里没有妇女，卧室里没有少女。同样，每个年龄、每个性别的人都从城里出来，仿佛迎接一位久已离家、从远方旅居归来的父亲。你们确实把他归还给了同道热忱者的团体。但天主的恩宠没有让他永远留在那里，而是再次将他移到河的对岸，以致乡间许多地方都充满了殉道者的馨香。即使当他来到这里，他也注定不会是孤独的，而是很快得到了一位邻居、一位同住者和一位同样生活的人。因为他与他分享了同样的尊位，并且为了宗教而表现出同样的胆量。因此，他获得了与殉道者相同的居所，这位了不起的人似乎不是殉道者空洞的模仿者。因为有很长一段时间他在那里劳苦，不断地给皇帝写信，使当局感到厌倦，并为殉道者带来身体的服事。因为你们无疑知道并记得，当夏日的正午太阳当空时，他常与他的相识们每天到那里去，不仅是作为旁观者，而且也是想要参与所进行的事。因为他常常搬动石头、拉绳索，并比工人本人更先听取想要建造房屋的人的意见。因为他知道，他知道这些事情为他存留着怎样的赏报。因此，他继续服务于殉道者，不仅是通过辉煌的建筑，也不是通过不断的盛宴，而是通过比这些更好的方法。这是什么？他效法他们的生活，仿效他们的勇气，始终按照他的能力，在自己内保持殉道者的形象鲜活。因为看，他们将自己的身体交于屠宰，他却抑制了自己在地上的肢体。他们止住了火焰，他却熄灭了情欲的火焰。他们与野兽的牙齿搏斗，而这个人却制服了我们最危险的情欲——愤怒。为这一切，让我们感谢天主，因为他如此赐给了我们高贵的殉道者，以及堪当殉道者的牧者，为成全圣徒，为建树基督的身体（见：\BibleRef{弗 4:12}），愿光荣、尊崇和权能归于圣父，偕同圣神，即赋予生命的圣神，从现在直到永远，世世无尽。阿们。\par

\SourceRecord{Translated by T.P. Brandram.；Nicene and Post-Nicene Fathers, First Series；Vol. 9.；Edited by Philip Schaff.；Buffalo, NY: Christian Literature Publishing Co.,；1889.；<http://www.newadvent.org/fathers/1906.htm>.}
